% =====================================================================
% Chapter 8: Discussion and Conclusions
%
% Required topic: the discussion of the research and the conclusions.
% The chapter closes the thesis with its conclusions and with the
% directions it leaves open.
%
% The conclusions are the ones of the Conclusion section of the paper.
% =====================================================================
\chapter{Conclusion}
\label{ch:discussion}



We present \tool, a verifier that accelerates constraint-solver-based
global robustness verification and is the first to support its
incremental verification. \tool tightens the encoding with
perturbation-difference intervals and zonotope bounds, relaxes neurons
with a novel conditional-triangle relaxation, and, in transfer mode,
reuses a prior verification through difference bounds and warm-start
hints. \tool speeds up \baseline by $7.6\times$ on
average, and by $12.4\times$ with transfer, on the instances
where at least one verifier terminates. On the
hardest tasks, where all verifiers time out, it returns bounds
$19.9\times$ tighter, and $20.9\times$
with transfer. Finally, our results show that the gains of \tool carry
over to larger networks, a network with
over $3\times$ more neurons than the largest network of \baseline and a
network with 500 neurons in one hidden layer, on which the bounds of
\baseline remain over $97\%$ of $\delta_{\max}$ apart at the timeout. We identify several future directions.

\section{Small Perturbations}
\label{sec:future-small-perturbations}
%
% TODO (Dana): Neta, could you shorten each of the three directions to one short paragraph: the problem, the idea, and what is left open? The longer they are, the more the examiners will ask about them in the defense, so keep only what you want to expand on there.


The initial intuition is that small perturbations would be easier to verify.
However, small perturbations can be very difficult to do so.
For examle, under an $\ell_\infty$ perturbation with a very
small $\ep$, every perturbed neuron sits within a hair of its clean
counterpart: the difference intervals
$[\Delta^{\hat z,\downarrow}_{m,k},\Delta^{\hat z,\uparrow}_{m,k}]$
are narrow, and the perturbed copy carries little that the clean one
does not already determine.

\tool profits from this far less than it might. Its per-neuron bounds
$[l_{m,k},u_{m,k}]$ are computed over the whole input domain, so which
neurons are unstable, and hence how many binaries the program carries,
is the same whether $\ep$ is large or vanishingly small. The
conditional triangle does not remove them either: relaxing $z_{m,k}$
conditions its triangle on the other copy's binary $a^{p}_{m,k}$, so
the perturbed copy keeps its binary exactly where the clean one gives
its up. The program still pays for two copies' worth of unstable
neurons in order to encode a second copy that has almost nothing left
to say. Moreover,  a small $\ep$
pins $\hat{\delta}_{c_t}$ near zero, so the hyper-adversarial attack
may not return a feasible pair, the search then starts with no lower
bound, and the solver is left to bring $\dup$ down on its own.

We can profit from the same closeness that causes the difficulty:
where \tool uses the difference intervals to couple the two copies, a
small enough perturbation makes them tight enough to stand in for
one. Here the perturbed copy is worth dropping altogether: we encode
the clean copy alone and replace each perturbed neuron by its
difference interval. That halves the binary variables the solver has to decide,
taking the square root of the worst-case search space, and gives up
only as much as the perturbation can move the logits, which in this
regime is little.

Formally, the difference intervals already carry the perturbation
along the network, from the input through each affine map and each
ReLU, so at the output they place every logit of the perturbed copy
against its clean counterpart,
\begin{equation}
  N(x)_k+\Delta^{\hat z,\downarrow}_{L,k}
  \;\leq\; N(x^{p})_k \;\leq\;
  N(x)_k+\Delta^{\hat z,\uparrow}_{L,k}.
  \label{eq:future-logit-difference}
\end{equation}
Now $\conf{N}{x^{p}}{c_t}>0$ says exactly that $c_t$ is the strict
argmax of $N(x^{p})$, so it is the conjunction over the other classes
of $N(x^{p})_{c_t}>N(x^{p})_k$, and therefore:
\[
  N(x)_{c_t}+\Delta^{\hat z,\uparrow}_{L,c_t}
  \;\geq\; N(x^{p})_{c_t}
  \;>\; N(x^{p})_k
  \;\geq\; N(x)_k+\Delta^{\hat z,\downarrow}_{L,k},
\]
whose two ends mention the clean copy alone. Rearranging them,
\begin{equation}
  N(x)_{c_t}-N(x)_k \;>\;
  \Delta^{\hat z,\downarrow}_{L,k}-\Delta^{\hat z,\uparrow}_{L,c_t}
  \qquad\text{for every } k\neq c_t.
  \label{eq:future-one-copy}
\end{equation}
Condition~\eqref{eq:future-one-copy} is necessary for
$\conf{N}{x^{p}}{c_t}>0$ but not sufficient. Encoding it in place of
the perturbed copy over-approximates the set of feasible pairs, so the
upper bound $\dup$ remains sound, while a feasible solution need not
correspond to a real pair $(x,x^{p})$. How much precision this costs
for a small $\ep$ is left for future research.

\section{Verifying New Perturbations in Global Robustness Settings}
\label{sec:future-perturbations}

Today, \tool, and the other verifiers we compare it to, support only a handful
of perturbations. However, there are many more perturbations that are of
interest to the research community, and we discuss below how \tool can be
extended to support them as well.

\paragraph{$\ell_0$ perturbation.}
An $\ell_0$ perturbation allows up to $t$ pixels to be modified within
$[0,1]$, leaving the rest untouched. Bounding it coordinate by
coordinate collapses the analysis: because any pixel could be one of
the $t$, interval propagation assigns $[-x_k, 1-x_k]$ to every
coordinate -- the very interval \tool already uses for a patch pixel
-- discarding the sparsity constraint entirely. We retain the budget
$t$ by evaluating the perturbation at the first affine layer rather
than at the pixel. Let $S=\{k \mid x^{p}_k \neq x_k\}$ be the pixels
the perturbation touches, so $\lvert S\rvert\leq t$. The bias cancels
in the difference of the two pre-activations, leaving a sum over $S$
alone,
\[
  \hat z^{p}_{1,j}-\hat z_{1,j}
  \;=\; \sum_{k} w_{1,j,k}\,(x^{p}_k-x_k)
  \;=\; \sum_{k\in S} w_{1,j,k}\,(x^{p}_k-x_k) .
\]
Each $\lvert x^{p}_k-x_k\rvert$ is at most $1$, as both lie in
$[0,1]$, so the difference is at most
$\sum_{k\in S}\lvert w_{1,j,k}\rvert$, and since the perturbation
chooses $S$ the bound must take its worst choice, the $t$ largest
weight magnitudes:
\begin{equation}
  \bigl\lvert \hat z^{p}_{1,j} - \hat z_{1,j} \bigr\rvert
  \;\leq\; \max_{\lvert S\rvert\leq t} \sum_{k\in S}\lvert w_{1,j,k}\rvert
  \;=\; \sum_{k \in T_j} \lvert w_{1,j,k} \rvert ,
  \label{eq:future-l0}
\end{equation}
where $T_j$ selects the $t$ coordinates with the largest weight
magnitudes $\lvert w_{1,j,k} \rvert$.

Encoding it is possible, but it can be expensive. Intuitively, we can
set an indicator per pixel, saying whether that pixel is perturbed or
not, which costs a binary for every input coordinate and so grows
with the size of the input: $784$ on MNIST and $3072$ on CIFAR-10,
against a program whose cost is already exponential in its binaries.
It also leaves $\binom{n}{t}$ choices of support to check. We can
instead rely on over-approximations such
as~\eqref{eq:future-l0}, or on the algorithms that already solve
$\ell_0$ for local robustness, where Calzone covers the subsets of
pixels with a covering design~\citep{shapira2023calzone} and CoVerD
chooses among such designs without building
them~\citep{shapira2024coverd}. Carrying either to a setting with two
copies and an unknown input requires further research.

\paragraph{$\ell_1$ perturbation.}
An $\ell_1$ perturbation enforces a global constraint across all coordinates:
$\sum_k \lvert x^{p}_k - x_k \rvert \leq \ep$. Because any single pixel may
consume the entire budget, coordinate-wise interval propagation assigns
$[-\ep, \ep]$ to every input coordinate. We can preserve the
global budget by evaluating the first affine layer via H\"older's inequality:
\[
  \bigl\lvert \hat z^{p}_{1,j} - \hat z_{1,j} \bigr\rvert
  \;\leq\; \ep \max_k \lvert w_{1,j,k} \rvert ,
\]
which replaces the sum over absolute weights with the single
largest row magnitude.

The budget here is cheap to encode: an auxiliary
$u_k$ per coordinate with $u_k \geq x^{p}_k - x_k$ and
$u_k \geq x_k - x^{p}_k$, together with the single row
$\sum_k u_k \leq \ep$, captures the $\ell_1$ ball exactly and adds no
binary. The difficulty lies in the bounds rather than in the
encoding, and it outlasts the first layer. The bound above spends the
budget once, at that layer, and from there on the propagation of
Section~\ref{sec:blend-pert-diff} is interval arithmetic over each
neuron separately, so two neurons may both sit at their worst case
although no single perturbation can place them there together: the
budget they share is gone. The zonotope cannot represent the
budget exactly either, as
it is a linear image of a box and the $\ell_1$ ball is not (for $n\geq 3$). Carrying a
shared budget through the network, rather than spending it at the
first layer, is what this direction needs and what we leave open.

\paragraph{$\ell_2$ perturbation.}
An $\ell_2$ perturbation enforces $\lVert x^{p} - x \rVert_2 \leq \ep$.
Coordinate-wise interval bounding relaxes this ball into an $\ell_\infty$
hypercube of radius $\ep$, admitting points with $\ell_2$-norm up to
$\sqrt{n}\ep$. By the Cauchy--Schwarz inequality, the deviation at the first
pre-activation is bounded directly by:
\[
  \bigl\lvert \hat z^{p}_{1,j} - \hat z_{1,j} \bigr\rvert
  \;\leq\; \ep \, \lVert w_{1,j,\cdot} \rVert_2 .
\]
All standard norms instantiate a single dual relation: for an $\ell_p$
perturbation, the pre-activation difference is bounded by
$\ep \lVert w_{1,j,\cdot} \rVert_q$ with $1/p + 1/q = 1$. The coordinate-wise
formulation in \tool corresponds to $q = 1$, which is tight only when
$p = \infty$. Downstream propagation (Section~\ref{sec:blend-pert-diff})
proceeds identically, requiring only this revised first-layer base case.
The single
quadratic row $\sum_k (x^{p}_k - x_k)^2 \leq \ep^2$ is convex, and can be
encoded as is to the solver.
The difficulty is again in the bounds, and it is the same difficulty.
Cauchy--Schwarz settles the first layer, and past it the propagation
treats each neuron's difference on its own, as though every one of
them could reach its worst case at once, which no perturbation of
norm $\ep$ can arrange. Here too the zonotope cannot represent the budget exactly, as the
$\ell_2$ ball is not a linear image of a box. The two norms therefore leave the same
question: how to carry one shared budget through a network whose
abstractions are built to carry a box. We regard this as the main
obstacle to supporting them and leave it to future research.

\section{Training Networks for Verifiability}
\label{sec:future-training}

Every direction above takes the network as given. The cost of the
analysis, however, is set by the network:
Section~\ref{sec:background-exact-verification} ties it to the number
of unstable neurons, and
Section~\ref{sec:future-small-perturbations} observes that this number
does not fall when the perturbation shrinks, because the bounds that
decide it are taken over the whole input domain. The network itself is
the remaining lever.

Training for verifiability is being studied for local robustness.
\citet{xiao2019relustability} add a regularizer penalizing
pre-activation bounds that straddle zero, which leaves more neurons
decided and shortens the analysis, and
\citet{chen2022grafting} replace unstable neurons with linear
functions whose slope and intercept are then retrained, carrying
complete verification to far larger networks.

To the best of our knowledge, this question has been asked only for
local robustness, never for exact verification of global robustness, and the two differ in
what they ask of a neuron.%
% TODO (Dana): GloRo nets~\citep{leino2021gloro}, cited in Related Work, train for a (Lipschitz-based) global robustness notion with built-in certification. Mention them here and say why they differ, or keep the qualifier.
 Our bounds are taken over all of
$[0,1]^n$, so a neuron is stable only if its pre-activation holds one
sign over the entire domain. That target is not merely harder but
bounded: were every
neuron stable over the whole domain, no ReLU would ever switch and the
network would be an affine map, so full stability reduces the classifier to a linear one. What to ask for is a frontier rather than a goal:
how many neurons can be decided over the whole domain before the
accuracy goes.

A second observation makes the target easier. \tool does not need a
neuron to be stable, only near-decided:
Section~\ref{sec:blend-condtri} relaxes it once the normalized
triangle area $\tfrac{u\,\lvert l\rvert}{2(u-l)}$ falls below $\tau$,
so a neuron whose pre-activation barely crosses zero is unstable and
yet costs no binary variable. Training against that area rather than against the sign asks
for less, is continuous where the sign is not, and optimizes the
quantity our relaxation actually tests. Whether it also buys back the
accuracy that whole-domain stability would cost is left open, as is
whether the objective should be the speed of the analysis at all,
rather than $\delta^*$ itself, which is what decides how many inputs a
deployed network has to flag.

The training regime deserves study in its own right. Adversarial
training, certified training and ordinary training leave a network in
different shapes, and we do not know how each of them bears on the
quantities above: on how many neurons stay decided over the whole
domain, on how small the triangle areas become, and on $\delta^*$
itself. Measuring that, rather than assuming it, is where we would
begin.
